\documentclass[11pt]{article}
\usepackage[margin=1in]{geometry}
\usepackage{amsmath,amssymb,amsthm}
\usepackage{enumitem}
\usepackage[colorlinks=true,allcolors=blue]{hyperref}

\newtheorem*{lemC}{Lemma 4 (FGL+19, Lemma~3.1 and its construction)}

\newcommand{\F}{\mathbb F}
\newcommand{\Gr}{\operatorname{Gr}_q}
\newcommand{\ZetaTransform}{\mathsf{Z}}
\newcommand{\ETH}{\mathrm{ETH}}
\DeclareMathOperator{\poly}{poly}

\title{What the Lean verification of\\ \emph{A Fast Algorithm for Subspace Convolution}
does not cover}
\author{Companion to the directory \texttt{lean/} of the repository
\texttt{shlomi-perles/subspace-convolution}}
\date{}

\begin{document}
\maketitle

\noindent
Every numbered result of the paper is formalized and checked by Lean using only the three standard
axioms:
\begin{itemize}
  \item Theorems~1, 2 and~3, Lemmas~1 and~2, Theorem~5, Corollary~2 and
    Theorem~7;
  \item Lemma~3, Theorem~6, Corollary~1, Proposition~1, Lemmas~5 and~6,
    Theorems~8, 9 and~10.
\end{itemize}
The unnumbered claims are formalized too: the estimates of Appendices~A and~B, the equations of
Sections~3--5, the applications of Section~5 and Appendix~C, and the footnote of Section~1.
This note lists, as precisely as possible, what is \emph{not} formalized: modelling choices,
hypotheses, and the few claims outside the model.

\section*{A. The transform part (Sections 1--5, Appendices A, B, C, D, F)}

\begin{enumerate}[leftmargin=1.6em]
\item \textbf{Arithmetic operations are gates of straight-line programs.}
  \begin{itemize}
    \item Linear computations are programs with gates $z=ax+by$; bilinear ones (span convolution,
      kernels whose coefficients are inputs, the count $\sum_U f(U)B_g(U)$) may also multiply two
      computed values.
    \item One gate is at most three ring operations, so the bounds are the paper's operation
      counts up to a factor of three.
    \item Programs are non-uniform: one program for each $n,k$ (and each coefficient ring).
  \end{itemize}
\item \textbf{Index construction is not formalized.} The paper's claims that constructing subspace
  bases, cosets, kernels and coordinate changes takes table size times $\poly(n)$ work
  (Section~1, Section~5.2 and the first paragraph of Appendix~D) concern a machine model. They are
  not formalized.
\item \textbf{Bit costs (Corollary~2).}
  \begin{itemize}
    \item Proved: the Möbius coefficients and the output bound; the choice of a modulus
      $L\in\{2B+1,2B+2\}$ coprime to $p$; exact recovery from centered residues; and the number of
      operations modulo $L$.
    \item Taken as the bit-cost model: each operation modulo $L$ is charged $(\log_2 L+1)^2$ bit
      operations (schoolbook arithmetic).
    \item Not formalized: the remarks after Corollary~2, namely
      \begin{itemize}
        \item evaluating over explicitly represented finite fields;
        \item the polynomial bit overhead of span convolution;
        \item determinant bounds for rational kernel coefficients and the resulting
          $\poly(n,H)$ bit costs of kernel applications and of solving.
      \end{itemize}
  \end{itemize}
\item \textbf{Theorem~5 outside the arithmetic model.}
  \begin{itemize}
    \item Formalized: both transforms (with their costs over $\mathbb Q$), the table $t$, and
      the identities $a=\ZetaTransform_n t$ and $C=\ZetaTransform_n^{-1}(2^{a})$.
    \item Accounted for in the documentation only: building $t$ ($\ell$ increments) and
      $2^0,\dots,2^\ell$ ($\ell$ doublings). Indexing the list by the integer $a(U)$ is a table
      lookup, not an arithmetic operation.
  \end{itemize}
\item \textbf{Invariant linear systems (Appendix~C).}
  \begin{itemize}
    \item Formalized: the structure of the small system~(eq. 30), the equivalence
      ``$K$ invertible $\iff$ $\det B\ne0$'', and the inverse kernel $C$ when $Bc=e_0$.
    \item Formalized: the cost of applying $C=K^{-1}$.
    \item Not counted as gates: solving the $(d+1)\times(d+1)$ system, which needs field
      divisions.
  \end{itemize}
\item \textbf{Statements about other algorithms are not formalized.}
  \begin{itemize}
    \item Running times of earlier algorithms (\cite{KKW16}, \cite{MOR04}, \cite{BTV10},
      \cite{GHO13}) and Kennes' addition-only lower bound \cite{Ken92} are cited facts. Formalized are only
      the counts they rest on (cover pairs, the inclusion count of Section~5.3).
    \item The eigenvalues of the Grassmann graph, which the paper cites from \cite{Sri14,MOR04},
      are \emph{proved} in Lean.
  \end{itemize}
\end{enumerate}

\section*{B. The coloring part (Section 6, Appendices G, H): hypotheses and unmodelled details}

\begin{enumerate}[leftmargin=1.6em]
\item \textbf{Lemma~4 is a hypothesis.} It is cited from Fomin--Golovach--Lokshtanov--Saurabh--Zehavi
  and enters Theorem~8 as a hypothesis about the cost model:
  \begin{lemC}
  Under $\ETH$, for any computable $f$, 4-Monotone min-CSP has no $f(k)n^{o(k)}$ algorithm, even
  when $|\mathcal C|=O(k^2n^2)$ and $L:=\max(1,\max_C|\mathcal R_C|)\le\max(n,4)$.
  \end{lemC}
\item \textbf{ETH and \#ETH are hypotheses.}
  \begin{itemize}
    \item Theorems~9 and~10 assume them in an abstract cost model.
    \item A Lean-checked refutation of ETH was published in October 2026
      (\texttt{github.com/anthropics/formal-math}, folder \texttt{eth/}); it has not been peer
      reviewed. If it is correct, the conditional forms of Theorems~8, 9 and~10 hold
      vacuously.
    \item The unconditional running-time transfers are proved, and they are unaffected.
  \end{itemize}
\item \textbf{Running times of the reductions are assumed.} The cost model assumes the following
  instead of proving them on a machine model:
  \begin{itemize}
    \item the constructions take $g_1(k)n^{O(1)}$ time (Theorem~8) and $2^{O(k)}\poly(m)$ time
      (Appendix~H);
    \item exact division by the known factor of eq.~(35) is cheap (Theorem~10).
  \end{itemize}
\item \textbf{Computability of $f$ is not tracked.}
  \begin{itemize}
    \item Theorem~8 is formalized for arbitrary real-valued $f$ on both sides.
    \item Applying Lemma~4 as stated needs the step ``the composite bound $F(k)$ is computable
      when $f$ is''. That step is not formalized.
  \end{itemize}
\item \textbf{Bit complexity of Theorem~6 is not modelled.}
  \begin{itemize}
    \item Theorem~6 and Corollary~1 are verified in arithmetic operations, counting the
      transitions of the aggregated dynamic program.
    \item Not formalized: the claims that all counts have $O(N\log N)$ bits, and that exact
      arithmetic costs a polynomial factor.
  \end{itemize}
\item \textbf{The cost of reconstructing an optimal coloring is not bounded.}
  \begin{itemize}
    \item The backtracking recursion and its correctness are formalized; its running time is not.
    \item Neither is the separate reconstruction for the linear-layout program of
      Corollary~1.
  \end{itemize}
\item \textbf{Unrooted rank-decompositions are not formalized.}
  \begin{itemize}
    \item A rank-decomposition is formalized as a rooted binary tree.
    \item The passage from a subcubic tree (``suppress degree-two vertices and subdivide an edge
      to root it'') is not formalized.
  \end{itemize}
\item \textbf{Appendix~H is verified for a variant grouping.}
  \begin{itemize}
    \item Lean makes the clause-local groups once per literal occurrence, so $g=2k+3m$. The paper
      makes them once per pair $(i,j\in J_i)$, so $g=2k+\sum_i|J_i|$.
    \item The counting identity, the vertex count and the width bound $2k+19$ are proved for this
      variant.
  \end{itemize}
\item \textbf{Rank-decomposition versions are only partly formalized.}
  \begin{itemize}
    \item For Theorems~8, 9 and~10 only the caterpillar width bound is proved: the
      caterpillar of the supplied layout has the same width.
    \item The corresponding cost-model statements for supplied rank-decompositions are not
      stated.
  \end{itemize}
\end{enumerate}

\begin{thebibliography}{9}
\bibitem{BTV10} B.-M.~Bui-Xuan, J.~A.~Telle, M.~Vatshelle, \emph{$H$-join decomposable graphs and
algorithms with runtime single exponential in rankwidth}, Discrete Appl.\ Math.\ 158(7), 2010.
\bibitem{FGL+19} F.~V.~Fomin, P.~A.~Golovach, D.~Lokshtanov, S.~Saurabh, M.~Zehavi,
\emph{Clique-width III: Hamiltonian Cycle and the Odd Case of Graph Coloring},
ACM Trans.\ Algorithms 15(1), 2019.
\bibitem{GHO13} R.~Ganian, P.~Hlin\v{e}n\'y, J.~Obdr\v{z}\'alek, \emph{A unified approach to
polynomial algorithms on graphs of bounded (bi-)rank-width}, European J.\ Combin.\ 34(3), 2013.
\bibitem{Ken92} R.~Kennes, \emph{Computational aspects of the M\"obius transformation of graphs},
IEEE Trans.\ Syst.\ Man Cybern.\ 22(2), 1992.
\bibitem{KKW16} P.~Kaski, J.~Kohonen, T.~Westerb\"ack, \emph{Fast M\"obius inversion in semimodular
lattices and ER-labelable posets}, Electron.\ J.\ Combin.\ 23(3), 2016.
\bibitem{MOR04} D.~K.~Maslen, M.~E.~Orrison, D.~N.~Rockmore, \emph{Computing isotypic projections
with the Lanczos iteration}, SIAM J.\ Matrix Anal.\ Appl.\ 25(3), 2004.
\bibitem{Sri14} M.~K.~Srinivasan, \emph{The Goldman--Rota identity and the Grassmann scheme},
Electron.\ J.\ Combin.\ 21(1), 2014.
\end{thebibliography}

\end{document}
